\section{\textcolor{Salmon}{Anforderungsanalyse anhand des Zustands des bestehenden Legacy-\-Systems und den Nutzerbedürfnissen} \textcolor{red}{Nutzerbedürfnisse? sicher?}}
\label{ch:Dashboard:sec:Anforderungen}
Überblick über Anforderungen an das bestehende System / das neue System

\noindent
Ableitung funktionaler und nicht-funktionaler Anforderungen\\
- Was muss das neue System mindestens können?\\
- was soll es besser oder zusätzlich leisten?\\

\noindent
- die Tabellen aus der ISO nutzen / ausfüllen \cite{DIN_EN_ISO_9241-110} \autoref{Add:ISO-Tabellen}\\
Erste Spalte: was ist auf so ein System anwendbar?\\
Zweite Spalte: was wurde auf das alte System angewendet?\\
Dritte Spalte: was wird sich vom neuen System erhofft?\\
Vierte Spalte: Was wurde im neuen System tatsächlich umgesetzt?

funktionale Anforderungen:
- Testergebnisse anzeigen
- Sammelpunkt für alle detailierteren Übersichtsseiten sein
- alle Ergebnisse auf einer Seite
- 

ISO-Tabellen:
- Anzahl mehr als altes System erfüllt

   alt neu Verschlechterung
6   2   5   0   Aufgabenangemessenheit
6   3   5   1   Selbstbeschreibungsfähigkeit
-> Informationen so darstellen, dass eindeutig ersichtlich ist, welche Benutzungsschnittstellen-Elemente interaktiv und welche nicht interaktiv sind - / - teilweise - Links / Buttons im Darkmode unscheinbar
8   5   6   1   Erwartungskonformität
-> Informationen auf unterschiedlichen Gerätearten und Anzeigegrößen angemessen darstellen - immer gleich, aber angemessen genug - das Grid ließ sich nicht fix einstellen
6   1   5   0   Erlernbarkeit
6   4   5   0   Steuerbarkeit
1   0   1   0   Robustheit gegen Benutzungsfehler
4   3  3+1  1   Benutzerbindung
-> Vertrauen schaffen hinsichtlich seiner Verwendung - zeigt Firmen Farben und Symbole - Firmenunabhängig